\documentclass[11pt]{article}
\usepackage[utf8]{inputenc}
\usepackage[T1]{fontenc}
\usepackage{amsmath,amssymb,amsfonts}
\usepackage{amsthm}

\newtheorem{theorem}{Theorem}[section]
\newtheorem{lemma}[theorem]{Lemma}

\usepackage{booktabs}
\usepackage{longtable}
\usepackage{array}
\usepackage{multirow}
\usepackage{float}
\usepackage{colortbl}
\usepackage{makecell}
\usepackage{xcolor}
\usepackage[margin=1in]{geometry}

\title{Minimum-Payment Expense Settlement}
\date{November 2026}

\begin{document}

\maketitle

\begin{abstract}
This document describes the mathematical foundations and algorithmic methods used in the Expenses Settler project, a Python library for computing the globally minimum number of payments required to settle shared expenses from a CSV input.

The core problem reduces to computing each participant's net balance (total paid minus total owed) and then finding a minimum-size multiset of transfers $(i, j, a)$ that zeroes all balances. We present three algorithms with different trade-offs between optimality, time, and space:

\begin{enumerate}
  \item An \textbf{exact} bitmask dynamic programming algorithm achieving the theoretical minimum $\min(n-c, \, n-1)$ payments in $O(3^n)$ time, where $n$ is the number of non-zero-balance people and $c$ is the maximum number of disjoint zero-sum subsets.
  \item A \textbf{heuristic} greedy subset-packing algorithm running in $O(n \cdot 2^n)$ time that may miss the optimum on crafted inputs.
  \item A \textbf{greedy} two-heap largest-debtor-to-largest-creditor algorithm running in $O(n \log n)$ time with no size limit but no optimality guarantee.
\end{enumerate}

We prove correctness for all three, establish the $n-c$ lower bound on payment count, and provide a complexity analysis that justifies the default threshold of $n=18$ for the exact method.
\end{abstract}

\section{Problem Formulation}

\subsection{Input}

A list of expenses, each represented as a triple $(p, A, U)$ where:
\begin{itemize}
  \item $p$ is the payer (a string identifier),
  \item $A$ is the total amount in non-negative integer cents,
  \item $U$ is the multiset of participants sharing the expense.
\end{itemize}

\subsection{Objective}

Given a collection of expenses, compute the \textbf{globally minimum} set of payments -- a list of triples $(i, j, a)$ where $a > 0$ is the amount sent from person $i$ to person $j$ -- such that after all payments, every person's net balance is zero.

\subsection{Notation}

\begin{itemize}
  \item $S$: the set of persons with non-zero net balance, with $|S| = n$.
  \item A \textbf{payment} is a triple $(i, j, a)$ with $a > 0$.
  \item A \textbf{settlement} is a multiset of payments $P$. Its size is $|P|$.
  \item A settlement is \textbf{feasible} if for every $i \in S$, the sum of amounts received minus the sum of amounts sent equals $-b_i$ (where $b_i$ is the net balance: positive = owed, negative = owes).
  \item A settlement is \textbf{minimal} if it is feasible and no feasible settlement uses fewer payments.
\end{itemize}

\section{Net Balance Computation}

\subsection{The Divmod Rounding Rule}

For an expense $(p, A, U)$ with $|U| = m$ participants, the share of each participant $u \in U$ is:

\[
  s_u = \begin{cases}
    \left\lfloor \frac{A}{m} \right\rfloor + 1 & \text{if } u \text{ is among the first } (A \bmod m) \text{ in CSV order}, \\
    \left\lfloor \frac{A}{m} \right\rfloor & \text{otherwise}.
  \end{cases}
\]

This implements \texttt{divmod(A, m)}: the first \texttt{(A mod m)} participants receive the extra cent each. The shares partition $A$ exactly:
\[
  \sum_{u \in U} s_u = m \cdot \left\lfloor \frac{A}{m} \right\rfloor + (A \bmod m) = A.
\]

\subsection{Net Balance Formula}

The net balance of person $i$ after all expenses is:
\[
  b_i = \sum_{e: \, p_e = i} A_e \;-\; \sum_{e: \, i \in U_e} s_{e,i}
\]

\begin{theorem}[Zero-Sum Invariant]
  $\sum_{i \in S} b_i = 0$.
\end{theorem}

\begin{proof}
  Every expense $A_e$ appears exactly once as a positive credit (to $p_e$) and exactly $|U_e|$ times with negative sign (as shares). Since the shares partition $A_e$ exactly, the total credit minus total debit for each expense is zero. Summing over all expenses gives zero.
\end{proof}

\subsection{Complexity}

$O(E \cdot m)$ where $E$ is the number of expenses and $m = \max |U_e|$.

\section{The Settlement Problem as a Graph Problem}

After computing $b_i$ for all $i \in S$, we form a directed graph where:

\begin{itemize}
  \item Debtors: $D = \{i \in S : b_i < 0\}$, with $d_i = -b_i$ (amount owed by $i$).
  \item Creditors: $C = \{i \in S : b_i > 0\}$, with $c_i = b_i$ (amount owed to $i$).
\end{itemize}

A payment $(i, j, a)$ reduces $d_i$ by $a$ and $c_j$ by $a$. A settlement is feasible iff all $d_i$ and $c_j$ reach zero.

\begin{theorem}[The $n-1$ Bound]
  Any feasible settlement uses at most $n-1$ payments.
\end{theorem}

\begin{proof}
  Sort persons by balance. Let person $p_1$ owe the most. Send $p_1$'s entire debt to any creditor. This eliminates at least one person. Repeat. At most $n-1$ payments needed.
\end{proof}

\begin{theorem}[The $n-c$ Bound]
  Let $c$ be the maximum number of pairwise disjoint zero-sum subsets of $S$. Then any settlement uses at least $n-c$ payments, and there exists a feasible settlement using exactly $n-c$ payments.
\end{theorem}

\begin{proof}
  \textbf{Lower bound.} Each zero-sum subset of size $k$ involves $k$ people. Eliminating all $k$ internal balances requires at least $k-1$ payments (same proof as the $n-1$ bound applied to the subset). Across $c$ disjoint subsets this accounts for $\sum (k_j - 1) = n - \sum k_j$ payments. The remaining $n - \sum k_j$ people need at most $(n - \sum k_j) - 1$ additional payments. Total $\geq n - c$.

  \textbf{Achievability.} For each zero-sum subset of size $k$, apply the greedy algorithm within the subset: exactly $k-1$ payments. Summing gives $n-c$ (or $n-1$ if $n-c > n-1$, i.e., $c < 1$).
\end{proof}

\section{The Exact Algorithm: Bitmask Dynamic Programming}

\subsection{High-Level Strategy}

The exact algorithm achieves the theoretical minimum $n-c$ by:
\begin{enumerate}
  \item Finding the \textbf{maximum} number $c$ of pairwise disjoint zero-sum subsets of $S$.
  \item For each subset, settling with greedy largest-debtor-to-largest-creditor (cost: $k-1$ for a subset of size $k$).
  \item Settling the remainder greedily (cost: $\leq (n - \sum k_j) - 1$).
\end{enumerate}

The optimal payment count equals $\min(n - c, \, n - 1)$.

\subsection{Step 1: Zero-Sum Mask Enumeration}

\textbf{Encoding.} Assign each person $p_j$ the index $j \in \{0, \ldots, n-1\}$. A subset $T \subseteq S$ is encoded as the bitmask $m \in \{0, 1, \ldots, 2^n-1\}$ where bit $j$ is 1 iff $p_j \in T$.

\textbf{Precomputation.} We precompute all zero-sum masks:
\[
  \texttt{zero\_sum\_masks} = \{m \mid \sum_{j: \text{bit } j \text{ of } m \text{ is set}} \!\! \texttt{balance\_list}[j] = 0\}
\]

\textbf{Complexity:} $O(n \cdot 2^n)$ time and space.

\textbf{Observation.} Since $\sum_{j=0}^{n-1} b_{p_j} = 0$ (Theorem 2), the full mask $111\ldots 1$ is always zero-sum.

\subsection{Step 2: Maximum Disjoint-Zero-Sum Count via DP}

\textbf{State.} Define:
\[
  \texttt{dp}[m] = \text{the maximum number of pairwise disjoint zero-sum subsets that can be found using only people whose bits are set in } m.
\]
The answer is $\texttt{dp}[2^n - 1]$.

\textbf{Base case.} $\texttt{dp}[0] = 0$.

\textbf{Transition.} For each $m > 0$, let $\texttt{lowbit} = m \;\&\; (-m)$ (the least significant set bit). We consider two cases:

\paragraph{Option A: Do not include lowbit's person in any subset.}
\[
  \texttt{dp}[m] = \texttt{dp}[m \oplus \texttt{lowbit}]
\]
This is always feasible: if the best partition of $m \oplus \texttt{lowbit}$ uses $c$ subsets, so does $m$ (we just ignore the lowbit person).

\paragraph{Option B: Include lowbit's person in a zero-sum subset.}
The subset must include lowbit (otherwise covered by Option A). Iterate over all submasks $s \subseteq (m \oplus \texttt{lowbit})$, forming $\texttt{candidate} = s \mid \texttt{lowbit}$:
\[
  \texttt{if candidate} \in \texttt{zero\_sum\_masks:}
\quad \texttt{dp}[m] = \max(\texttt{dp}[m],\; \texttt{dp}[m \oplus \texttt{candidate}] + 1)
\]

\textbf{Key insight.} The person at lowbit \textbf{must} belong to some subset (if they belong to one at all). We iterate only over zero-sum subsets that \textbf{include} lowbit, avoiding redundant sub-problems and ensuring every partition is considered exactly once.

\textbf{Complexity of the transition.} The number of submasks of a mask of weight $w$ is $2^{w-1}$. Summing over all $2^n$ masks gives:
\[
  \sum_{w=1}^{n} \binom{n}{w} \cdot 2^{w-1} = \frac{1}{2}\sum_{w=0}^{n}\binom{n}{w}2^w = \frac{3^n}{2} = O(3^n)
\]

\subsection{Step 3: Reconstruction via Parent Pointers}

Maintain $\texttt{parent}[m]$ = the mask of the last zero-sum subset added in the optimal solution:
\[
  \texttt{parent}[m] = \begin{cases}
    0 & \text{if Option A is chosen},\\
    \texttt{candidate} & \text{otherwise.}
  \end{cases}
\]

\textbf{Reconstruction algorithm:}
\begin{verbatim}
selected = []
m = (1 << n) - 1
while m > 0:
    sub = parent[m]
    if sub == 0: break          # no more subsets to extract
    selected.append(people set from sub)
    m = m ^ sub                  # remove extracted subset
return selected
\end{verbatim}

\begin{lemma}[Correctness of reconstruction]
The sequence of subsets produced by repeatedly extracting $\texttt{parent}[m]$ is a valid partition: extracted subsets are pairwise disjoint, all are zero-sum, and the union covers the subset of $m$ for which $\texttt{parent}$ was set.
\end{lemma}

\begin{proof}
By induction on the number of extraction steps. Each iteration removes exactly $\texttt{parent}[m]$ from $m$, and $\texttt{parent}[m]$ is a zero-sum mask (verified during DP). Since $m$ decreases strictly, the process terminates. The final $m$ with $\texttt{parent} = 0$ has no subsets remaining.
\end{proof}

\subsection{Step 4: Intra-Subset Settlement}

For each discovered zero-sum subset $T$ of size $k$, apply the greedy largest-debtor-to-largest-creditor within $T$ only. Since the subset is zero-sum, this terminates with all internal balances at zero in exactly $k-1$ payments.

\subsection{Step 5: Remainder Settlement}

People not in any extracted subset are settled by greedy pairwise matching over the entire balance set.

\subsection{Total Payment Count}

\[
  |P| = (k_1-1) + (k_2-1) + \cdots + (k_c-1) + |P_{\text{remainder}}|
\]
where $\sum k_j = \sum |T_j|$. If no remainder people exist: $|P| = n - c$. Otherwise $|P| \leq n - c - 1$. In all cases $|P| \leq n-1$ (Theorem 3) and $|P| \geq n-c$ (Theorem 4).

\subsection{Correctness}

\begin{theorem}[Optimality of the Exact Algorithm]
The exact algorithm produces a settlement of minimum size.
\end{theorem}

\begin{proof}
By Lemma 1, the DP finds the true maximum $c$. By Theorem 4, $n-c$ is a lower bound on any feasible settlement, and we construct one of that size. Therefore the result is optimal.
\end{proof}

\subsection{Complexity Summary}

\begin{table}[H]
\centering
\begin{tabular}{lll}
\toprule
Step & Time & Space \\
\midrule
Zero-sum mask enumeration & $O(n \cdot 2^n)$ & $O(2^n)$ \\
DP computation & $O(3^n)$ & $O(2^n)$ \\
Reconstruction & $O(n \cdot c)$ & $O(c)$ \\
Greedy intra-subset & $O(n \cdot c)$ & $O(n)$ \\
Greedy remainder & $O(n \log n)$ & $O(n)$ \\
\midrule
\textbf{Total} & $\mathbf{O(3^n)}$ & $\mathbf{O(2^n)}$ \\
\bottomrule
\end{tabular}
\caption{Exact algorithm complexity analysis.}
\end{table}

The exponential bound explains the default threshold of $n = 18$: $3^{18} \approx 3.87 \times 10^8$ operations is feasible in seconds, while $3^{24} \approx 2.82 \times 10^{11}$ is infeasible.

\section{The Heuristic Algorithm}

\subsection{Strategy}

Enumerate all zero-sum subsets in order of increasing size, then greedily select non-overlapping ones:
\begin{verbatim}
for size = 2 to n:
    for each combination of size persons:
        if sum of their balances is zero:
            add to zero_sum_subsets

sort zero_sum_subsets by size (smallest first)
selected = []; used = empty set
for subset in zero_sum_subsets:
    if subset \\cap used == \\emptyset:
        selected.append(subset)
        used |= subset
\end{verbatim}

\subsection{Complexity}

\begin{itemize}
  \item Enumeration: $\sum_{k=2}^{n} \binom{n}{k} = 2^n - n - 1$ subsets. Each check costs $O(n)$. Total: $O(n \cdot 2^n)$.
  \item Sorting: $O(n^2 \cdot 2^n)$.
  \item Greedy selection: $O(n \cdot 2^n)$.
  \item \textbf{Dominant term:} $O(n \cdot 2^n)$.
\end{itemize}

\subsection{Not Guaranteed Optimal}

The heuristic may select a large subset that blocks two smaller disjoint subsets. For example, with balances chosen such that a size-3 subset consumes people needed for two disjoint size-2 subsets, the heuristic produces more total payments than the exact DP.

\section{The Greedy Algorithm: $O(n \log n)$}

\subsection{Strategy}

Settle all balances in a single greedy pass using two min-heaps:

\begin{itemize}
  \item \textbf{Debtors heap:} keyed by $(-b_i, i)$, so the smallest key = the largest debtor.
  \item \textbf{Creditors heap:} keyed by $(-b_i, i)$, so the smallest key = the largest creditor.
\end{itemize}

\textbf{Why negate.} Python's \texttt{heapq} is a min-heap. Negating the balance value makes the smallest heap entry correspond to the largest absolute balance.

\subsection{Algorithm}

\begin{verbatim}
debtors = minheap[(-b_i, i) for i if b_i < 0]
creditors = minheap[(-b_i, i) for i if b_i > 0]

while debtors and creditors:
    k_d, d = debtors[0]      # k_d = -b_d < 0
    k_c, c = creditors[0]    # k_c = -b_c < 0
    amount = min(-k_d, -k_c)
    emit payment(d, c, amount)
    pop both heaps
    push back d if balance not exhausted
    push back c if balance not exhausted
\end{verbatim}

\textbf{Termination.} Each iteration reduces the sum of remaining balances by $\texttt{amount} > 0$. The sum is finite and non-negative, so the process terminates.

\textbf{Correctness.} At termination, at most one heap is non-empty. If both are non-empty, the heap-invariant (smallest = largest) implies a positive balance exists in both, so another payment can be made -- contradiction. Exactly one heap is non-empty. But since $\sum b_i = 0$ (Theorem 2), the remaining balances must all be zero.

\subsection{Payment Count and Gaps}

Each iteration that fully clears both a debtor and a creditor removes one person from consideration (one payment per removal). The upper bound is $n-1$ payments (at least one person removed per payment).

\textbf{Not guaranteed minimal.} In multi-group scenarios, the greedy may bridge zero-sum groups that could be settled more efficiently in isolation. The gap between greedy and optimal can be arbitrarily large for sufficiently structured inputs (multiple disjoint zero-sum components connected through a single "bridge" person).

\subsection{Complexity}

\begin{itemize}
  \item Heap construction: $O(n)$.
  \item Each iteration pops one person and pushes back at most one. Total iterations $\leq 2n-1$.
  \item Each push/pop: $O(\log n)$.
  \item \textbf{Total:} $O(n \log n)$ time, $O(n)$ space.
\end{itemize}

\section{Comparison Summary}

\begin{table}[H]
\centering
\begin{tabular}{llll}
\toprule
Property & Exact & Heuristic & Greedy \\
\midrule
Optimality & $\checkmark$ provably optimal & $\times$ not guaranteed & $\times$ not guaranteed \\
Time & $O(3^n)$ & $O(n \cdot 2^n)$ & $O(n \log n)$ \\
Space & $O(2^n)$ & $O(n \cdot 2^n)$ & $O(n)$ \\
Size limit & $n \leq 18$ (default) & any $n$ & any $n$ \\
Method & Bitmask DP & Greedy subset packing & Two-heap pairwise \\
\bottomrule
\end{tabular}
\caption{Algorithm comparison.}
\end{table}

\section{Edge Cases and Invariants}

\subsection{Zero-Value Expenses}
An expense with $A = 0$ contributes $0$ to every participant's share. No effect on balances.

\subsection{Single Participant}
An expense with $|U| = 1$ credits the payer $A$ and debits the participant $A$ by $A$. Net: $b_{\text{payer}} \leftarrow b_{\text{payer}} + 0$. No effect.

\subsection{Payer Not a Participant}
The payer receives $+A$ (credited) but has no negative share. Net: $b_{\text{payer}} \leftarrow b_{\text{payer}} + A$. The payer is a creditor for that expense.

\subsection{Payer Is a Participant}
The payer receives $+A$ and a negative share $s_p$. Net: $b_p \leftarrow b_p + A - s_p$. Since $s_p \leq A/m < A$ when $|U| \geq 2$, the payer always has a positive net from their own expense.

\subsection{Invariants}

\begin{enumerate}
  \item \textbf{All payment amounts are positive.} Every payment amount is $\min(-d_{\text{debtor}}, c_{\text{creditor}})$ with both arguments strictly positive (we only consider non-zero balances).

  \item \textbf{No self-payment.} The debtor heap and creditor heap are disjoint sets (balance $< 0$ vs. balance $> 0$). No person can appear in both.
\end{enumerate}

\section{Source Reference}

All algorithms are implemented in \texttt{settle\_expenses.py}:

\begin{center}
\begin{tabular}{lll}
\toprule
Function & Lines & Role \\
\midrule
\texttt{net\_balances()} & \S2 & Compute per-person net balance from expenses \\
\texttt{settle\_expenses()} & \S3 & Main entry; dispatches to exact/heuristic/greedy \\
\texttt{\_settle()} & \S4 & Greedy largest-debtor-to-largest-creditor within a group \\
\texttt{\_greedy\_settle\_balances()} & \S6 & Polynomial $O(n \log n)$ two-heap settlement \\
\texttt{\_find\_max\_zero\_sum\_subsets\_exact()} & \S4 & Bitmask DP, optimal subset count \\
\texttt{\_find\_max\_zero\_sum\_subsets\_heuristic()} & \S5 & Greedy subset enumeration, non-optimal \\
\bottomrule
\end{tabular}
\end{center}

\section{Dispatch Logic}

The \texttt{threshold} parameter of \texttt{settle\_expenses()} controls when "auto" falls back:

\[
  \text{effective method} = \begin{cases}
    \texttt{exact} & \text{if } \texttt{method} = \texttt{"exact"} \\
    \texttt{greedy} & \text{if } \texttt{method} = \texttt{"greedy"} \\
    \texttt{greedy} & \text{if } \texttt{method} = \texttt{"auto"} \text{ and } n > \texttt{threshold} \\
    \texttt{exact} & \text{if } \texttt{method} = \texttt{"auto"} \text{ and } n \leq \texttt{threshold}
  \end{cases}
\]

With default \texttt{threshold} $= 18$, the exact DP is used up to 18 people, and a \texttt{UserWarning} is emitted when falling back.

\end{document}
